% Options for packages loaded elsewhere
\PassOptionsToPackage{unicode}{hyperref}
\PassOptionsToPackage{hyphens}{url}
\documentclass[
  ignorenonframetext,
]{beamer}
\newif\ifbibliography
\usepackage{pgfpages}
\setbeamertemplate{caption}[numbered]
\setbeamertemplate{caption label separator}{: }
\setbeamercolor{caption name}{fg=normal text.fg}
\beamertemplatenavigationsymbolsempty
% remove section numbering
\setbeamertemplate{part page}{
  \centering
  \begin{beamercolorbox}[sep=16pt,center]{part title}
    \usebeamerfont{part title}\insertpart\par
  \end{beamercolorbox}
}
\setbeamertemplate{section page}{
  \centering
  \begin{beamercolorbox}[sep=12pt,center]{section title}
    \usebeamerfont{section title}\insertsection\par
  \end{beamercolorbox}
}
\setbeamertemplate{subsection page}{
  \centering
  \begin{beamercolorbox}[sep=8pt,center]{subsection title}
    \usebeamerfont{subsection title}\insertsubsection\par
  \end{beamercolorbox}
}
% Prevent slide breaks in the middle of a paragraph
\widowpenalties 1 10000
\raggedbottom
\AtBeginPart{
  \frame{\partpage}
}
\AtBeginSection{
  \ifbibliography
  \else
    \frame{\sectionpage}
  \fi
}
\AtBeginSubsection{
  \frame{\subsectionpage}
}
\usepackage{iftex}
\ifPDFTeX
  \usepackage[T1]{fontenc}
  \usepackage[utf8]{inputenc}
  \usepackage{textcomp} % provide euro and other symbols
\else % if luatex or xetex
  \usepackage{unicode-math} % this also loads fontspec
  \defaultfontfeatures{Scale=MatchLowercase}
  \defaultfontfeatures[\rmfamily]{Ligatures=TeX,Scale=1}
\fi
\usepackage{lmodern}
\ifPDFTeX\else
  % xetex/luatex font selection
\fi
% Use upquote if available, for straight quotes in verbatim environments
\IfFileExists{upquote.sty}{\usepackage{upquote}}{}
\IfFileExists{microtype.sty}{% use microtype if available
  \usepackage[]{microtype}
  \UseMicrotypeSet[protrusion]{basicmath} % disable protrusion for tt fonts
}{}
\makeatletter
\@ifundefined{KOMAClassName}{% if non-KOMA class
  \IfFileExists{parskip.sty}{%
    \usepackage{parskip}
  }{% else
    \setlength{\parindent}{0pt}
    \setlength{\parskip}{6pt plus 2pt minus 1pt}}
}{% if KOMA class
  \KOMAoptions{parskip=half}}
\makeatother
\usepackage{longtable,booktabs,array}
\usepackage{caption}
\captionsetup[table]{skip=6pt}
\newcounter{none} % for unnumbered tables
\usepackage{calc} % for calculating minipage widths
\usepackage{caption}
% Make caption package work with longtable
\makeatletter
\def\fnum@table{\tablename~\thetable}
\makeatother
\setlength{\emergencystretch}{3em} % prevent overfull lines
\providecommand{\tightlist}{%
  \setlength{\itemsep}{0pt}\setlength{\parskip}{0pt}}
% Manuscript macro/package fallbacks (newcommand -> providecommand).
\usepackage{amsmath}
\usepackage{amssymb}
\usepackage{booktabs}
% Contents listing: article.cls prefixes every \l@section entry with
% \addvspace{1.0em}, which across ten sections costs about ten lines and pushed
% the ~40-entry listing one line past page 2 — leaving a third sheet carrying a
% single "10 References" line. Tighten that inter-section gap for the duration
% of \tableofcontents only, inside a group, so body spacing is untouched and no
% entry is dropped (reducing \tocdepth would have hidden 29 real subsections).
  \providecommand{\addvspace}[1]{\vskip 2\p@}%
% Document metadata. config.yaml declares a subtitle and a keyword list, and
% the producer emits both as tokens, but the template's \hypersetup writes only
% pdftitle/pdfauthor/pdflang -- so `pdfinfo` reported no Subject and no
% Keywords. This block runs after the template's, so it wins.
%
% The two values below are WRITTEN by scripts/z_generate_manuscript_variables.py
% (sync_preamble_metadata), never typed. A double-brace token cannot be used here:
% preamble.md is substituted into output/manuscript/ like any other section, but
% the template's _manuscript_source.py then copies the raw file back over that
% copy, so an unresolved token would reach hyperref and land in the PDF's
% metadata verbatim.
% Bounded identifier splitting. The template defines
%   \protected\def\breaktt#1{\begingroup\ttfamily\seqsplit{#1}\endgroup}
% and \seqsplit permits a break after EVERY character with no continuation cue.
% In the 2026-09-05 PDF that rendered `model_family_manifest.json` as
% `model_family_manifest.js` + `on` and `src/manuscript_variables.py` as
% `(src/manusc` + `ript_variables.py` — the first is itself a valid filename, so
% a reader cannot tell a line break from a name. Keep seqsplit's guarantee that
% no code token overflows the measure, but forbid breaks inside the last
% \GNNttTail characters (an extension is never orphaned) and inside the first
% \GNNttHead (a one- or two-letter stub is never left behind).

% Slide layout defaults for warning-clean scientific decks.
\usepackage{etoolbox}
\IfFileExists{xurl.sty}{\usepackage{xurl}}{}
\IfFileExists{seqsplit.sty}{\usepackage{seqsplit}}{\newcommand{\seqsplit}[1]{#1}}
\protected\def\breakseq#1{\seqsplit{#1}}
\protected\def\breaktt#1{\begingroup\ttfamily\seqsplit{#1}\endgroup}
\setlength{\emergencystretch}{6em}
\tolerance=5000
\hbadness=10000
\hfuzz=1pt
\setlength{\tabcolsep}{2pt}
\AtBeginEnvironment{longtable}{\tiny\renewcommand{\arraystretch}{0.86}\setlength{\tabcolsep}{1pt}}
\AtBeginEnvironment{tabular}{\tiny\renewcommand{\arraystretch}{0.86}\setlength{\tabcolsep}{1pt}}
\AtBeginEnvironment{equation}{\tiny}
\AtBeginEnvironment{equation*}{\tiny}
\AtBeginEnvironment{align}{\tiny}
\AtBeginEnvironment{align*}{\tiny}
\AtBeginEnvironment{itemize}{\footnotesize}
\AtBeginEnvironment{enumerate}{\footnotesize}
\AtBeginEnvironment{description}{\footnotesize}
\setbeamerfont{caption}{size=\tiny}
\setbeamerfont{caption name}{size=\tiny}
\setbeamerfont{normal text}{size=\small}
\setbeamerfont{frametitle}{size=\small}
\setbeamerfont{section title}{size=\footnotesize}
\setbeamerfont{subsection title}{size=\footnotesize}
\setbeamertemplate{section page}{%
  \centering
  \begin{beamercolorbox}[sep=12pt,center,wd=\paperwidth]{section title}
    \parbox{0.86\paperwidth}{\centering\usebeamerfont{section title}\insertsection\par}
  \end{beamercolorbox}
}
\setbeamertemplate{subsection page}{%
  \centering
  \begin{beamercolorbox}[sep=8pt,center,wd=\paperwidth]{subsection title}
    \parbox{0.86\paperwidth}{\centering\usebeamerfont{subsection title}\insertsubsection\par}
  \end{beamercolorbox}
}
\setlength{\abovecaptionskip}{2pt}
\setlength{\belowcaptionskip}{0pt}

% Natbib and cross-reference fallbacks — slides don't load natbib
% or cleveref, but combined-PDF manuscript prose may emit these
% commands. The fallback renders citations as a bracketed key list
% and cross-references as detokenized labels so slides don't fail on
% undefined control sequences or raw underscores. \providecommand is
% a no-op if packages load later.
\providecommand{\citep}[1]{[#1]}
\providecommand{\citet}[1]{#1}
\providecommand{\citealp}[1]{#1}
\providecommand{\citeauthor}[1]{#1}
\providecommand{\citeyear}[1]{#1}
\providecommand{\cref}[1]{\texttt{\detokenize{#1}}}
\providecommand{\Cref}[1]{\texttt{\detokenize{#1}}}
\providecommand{\crefrange}[2]{\texttt{\detokenize{#1}}--\texttt{\detokenize{#2}}}
\providecommand{\Crefrange}[2]{\texttt{\detokenize{#1}}--\texttt{\detokenize{#2}}}
\providecommand{\paragraph}[1]{\textbf{#1}\ }

\newtheorem{proposition}{Proposition}
\newtheorem{hypothesis}{Hypothesis}
\newtheorem{remark}[theorem]{Remark}
\newtheorem{axiom}{Axiom}
\newtheorem{property}{Property}
\usepackage{bookmark}
\IfFileExists{xurl.sty}{\usepackage{xurl}}{} % add URL line breaks if available
\urlstyle{same}
% fallback for those not using the hyperref driver hyperxmp:
\makeatletter
\@ifundefined{xmpquote}{\newcommand{\xmpquote}[1]{#1}}{}
\makeatother
\hypersetup{
  hidelinks,
  pdfcreator={LaTeX via pandoc}}

\author{\texorpdfstring{}{}}
\date{}

\begin{document}

\section{Limitations and Next Steps}\label{sec:limitations_next_steps}

\begin{frame}[fragile,allowframebreaks]{Current Limitations}
\protect\phantomsection\label{current-limitations}
GNN 3.2.0 is a working system with deliberately scoped boundaries, and
it is worth stating those boundaries plainly rather than implying
broader completeness than the artifacts support. The reliability gates
that underpin our confidence in the system --- the semantic-fidelity
ledger and the cross-framework reliability ledger --- are reproducible
by command and recorded against the 9 model families, but this
manuscript does not assert a particular full-suite pass rate as a
headline number. Pass and skip counts shift with the local toolchain,
the presence or absence of optional simulation dependencies, and whether
the optional Ollama integration paths are exercised, so we treat the
reproducing command as the durable claim and leave the run-specific
tallies to the ledgers and execution traces that the pipeline writes
alongside each run. A reader who wants a number should regenerate it in
\end{frame}

\begin{frame}[fragile,allowframebreaks]{Current Limitations}
their own environment rather than trust a transcribed digit here.

A second limitation lives at the boundary between the model families and
the rendering backends. GNN registers 9 simulation backends, but
coverage across the family-by-backend matrix is intentionally uneven,
and the gaps are recorded as explicit \emph{profiled-unsupported}
statuses rather than silently omitted. The cross-framework acceptance
check is anchored on the continuous family, where JAX, NumPyro,
RxInfer.jl are compared directly; the continuous and hierarchical
families, by contrast, carry profiled-unsupported markers at the render
and execute steps (Steps 11 and 12) for backends that cannot yet
faithfully realize their continuous-state or temporal-depth structure.
This is honest bookkeeping, not a defect to be hidden: an unsupported
status with a recorded reason is more trustworthy than a forced
translation that quietly misrepresents the model. It also reflects the
\end{frame}

\begin{frame}[fragile,allowframebreaks]{Current Limitations}
state of the surrounding literature, where expected-free-energy
formulations, variational-inference reductions, and hierarchical
active-inference methods are still being actively refined rather than
collapsed into one settled executable recipe (Champion et al. 2024;
Nuijten et al. 2026; Rangarajan and Rao 2026). The practical consequence
is that not every model expressible in the GNN syntax can be executed on
every registered backend today, and authors should consult the
profiled-unsupported ledger that
\breaktt{scripts/run\_model\_family\_acceptance.py} writes --- the
command is given in sec.~6 --- before assuming a
given family will round-trip through a given framework.

\end{frame}

\begin{frame}[fragile,allowframebreaks]{Current Limitations}
A third limitation is one of scope rather than capability. The formal
account in sec.~3 fixes the generative-model
components \texttt{A}, \texttt{B}, \texttt{C}, \texttt{D} and \texttt{E}
as given and defines only state inference (eq.~6) and policy
inference (eq.~7). GNN can also express \emph{parameter}
learning: a matrix may be declared as a latent variable with a Dirichlet
prior, and the repository ships both a corpus model that does so
(\breaktt{input/gnn\_files/learning/dirichlet\_likelihood\_learning.md},
which learns its likelihood matrix \texttt{A} from observations by joint
variational message passing) and the renderer strategy that emits code
for it (\breaktt{src/render/rxinfer/\_strategies\_learning.py}). That
corpus directory is registered in no manifest family --- the
unregistered set is exactly \breaktt{input/gnn\_files/learning} --- so
parameter learning is exercised by neither the semantic-fidelity nor the
\end{frame}

\begin{frame}[fragile,allowframebreaks]{Current Limitations}
cross-framework ledger and carries no figure or table in this
manuscript. We state it here as implemented-but-out-of-scope rather than
leave the omission to be inferred from the family tables.

Finally, several capabilities depend on optional dependencies that are
not part of the minimal install. The simulation frameworks themselves,
the audio sonification path, the LLM-enhanced analysis step, and the
interactive visualization step all require extra packages that a lean
checkout will not have, and the corresponding pipeline steps degrade to
recorded skips rather than failures when those dependencies are absent.
This keeps the core parse-validate-export-visualize spine runnable in
constrained environments, but it means a full 0--24 traversal of all 25
steps reflects the optional surface that a particular machine has
installed. We consider this an acceptable engineering trade for
portability, but it is a limitation a reader should hold in mind when
\end{frame}

\begin{frame}[fragile,allowframebreaks]{Current Limitations}
interpreting any single end-to-end run.
\end{frame}

\begin{frame}[allowframebreaks]{Next Steps}
\protect\phantomsection\label{next-steps}
The roadmap is concrete and staged. The current release is GNN 3.2.0
(``Exemplar Gold Standard'', 2026-09-02). It builds on the three
safe-by-design orchestration contracts introduced in GNN 3.0.0, which
generate and validate data only, with no live infrastructure mutation.
The first is \emph{durable observation streams}: file- and array-backed
stream manifests (content-checksummed) with replayable execution traces,
so that an extended run can be observed, paused, and re-derived
deterministically before any live sensor or device-backed stream is ever
introduced. The second is \emph{resumable run sessions}: run-session
manifests with atomic checkpoint/resume, status inspection, and
cancellation-safe cleanup so that extended model-family acceptance runs
can be interrupted and resumed without corrupting partial state. The
third is \emph{auditable container plans}: generating hardened container
\end{frame}

\begin{frame}[allowframebreaks]{Next Steps}
plans with an explicit static security-review and rollback descriptor,
deliberately stopping short of mutating any real cluster. Each contract
is covered by tests over real objects with negative controls and a
strict acceptance gate, additive live wiring connects them to the
session-acceptance and run-manifest paths, and three new Model Context
Protocol tools expose them. The unifying discipline across all three is
that orchestration becomes more capable only as fast as it becomes more
inspectable.

\end{frame}

\begin{frame}[allowframebreaks]{Next Steps}
The next major target, v4.0.0, pushes toward bounded autonomy, and it is
gated behind the 3.0.0 orchestration work for good reason. The intent is
to promote today's proposal-only candidate-scoring machinery --- which
already ranks candidate patches using the existing validators,
model-family ledgers, and interpretability reports --- toward
\emph{reviewed self-editing of GNN files}. Crucially, the design keeps a
human in the loop: edits are proposed and applied only after explicit
user approval, with no automatic source mutation, and the autonomy is
wrapped in policy, rollback, and audit controls before any
self-modifying or distributed action is permitted. This continues the
same principle that governs the rest of the system, in the lineage of
active-inference accounts of self-organizing systems that act to
minimize free energy under explicit generative models (Friston 2010;
Parr et al. 2022): an agent should expand its license to act only in
\end{frame}

\begin{frame}[allowframebreaks]{Next Steps}
proportion to the reviewability of what it does. The v4.0.0 items are
not claimed as complete here; they are the recorded direction of travel,
and each will be evidenced by its own gates when it lands, just as the
3.0.0 contracts are evidenced by theirs today.
\end{frame}

\end{document}
